\documentclass[11pt]{article}
\usepackage[margin=1in]{geometry}
\usepackage{amsmath, amssymb}
\usepackage{booktabs}
\usepackage{listings}
\usepackage{xcolor}
\usepackage{hyperref}
\usepackage{longtable}
\lstset{
  basicstyle=\ttfamily\small,
  breaklines=true,
  frame=single,
  columns=fullflexible,
}

\title{Independent Replication of\\
       \emph{Efficient networks for quantum factoring}\\
       (Beckman, Chari, Devabhaktuni, Preskill; arXiv:quant-ph/9602016)}
\author{QC-200 Replication --- Ollie subagent \\ 2026-07-05}
\date{2026-07-05}

\begin{document}
\maketitle

\begin{abstract}
We reproduce the concrete Sec.~VII (``$N=15$'') construction of Beckman,
Chari, Devabhaktuni, and Preskill --- their explicit low-gate network for the
modular-exponentiation operator $\mathrm{EXP\,N}(x{=}7, N{=}15)$ (Eq.~7.5) used
for the ``proof-of-principle'' factoring-of-15 demonstration on a
Cirac--Zoller ion trap. Running the exact circuit in \texttt{qiskit 2.5.0} with
\texttt{qiskit\_aer 0.17.2}, we verify (i) that the operator computes the
Eq.~(7.3) lookup table $|a\rangle|0\rangle \mapsto |a\rangle|7^a \bmod 15\rangle$
correctly on all four inputs $a \in \{0,1,2,3\}$ with unit probability, (ii)
that the entangled state $\frac{1}{2}\sum_a |a\rangle|7^a \bmod 15\rangle$ is
prepared with fidelity $1.000000$ against the paper's Eq.~(7.2), (iii) that
the gate composition is exactly $[6, 0, 4]$ (Eq.~7.6), giving a Cirac--Zoller
cost of \textbf{34 laser pulses}, matching the paper exactly, and (iv) that
the $L=2$ QFT applied to the joint state gives $y$ uniformly on $\{0,1,2,3\}$
as predicted by Eq.~(7.10). \textbf{Verdict: REPLICATED} for the Sec.~VII
$N=15$ construction (headline claim); the abstract's ``38 pulses'' number
requires the further optimized $\mathrm{EXP\,N}'$ variant (Eq.~7.9), which we
identify but do not separately re-count (the 34-pulse first construction
plus 6-pulse QFT + 2-pulse Hadamard prep gives 42, matching the paper's
own intermediate number).
\end{abstract}

\section{Paper summary}

\textbf{Authors (verified from the fetched PDF's title page):} David Beckman,
Amalavoyal N.\ Chari, Srikrishna Devabhaktuni, John Preskill (Caltech;
CALT-68-2021). \textbf{Submitted:} 1996-02-21 to arXiv. \textbf{Topic:} explicit
gate-level networks for Shor's factoring algorithm on a Cirac--Zoller linear
ion trap, minimizing memory qubits and elementary operation counts.

Their central general-purpose bookkeeping result is that a $K$-bit number
can be factored in time $\mathcal{O}(K^3)$ using $5K+1$ qubits, with the
modular-exponentiation subcircuit needing about $72 K^3$ elementary gates
(equivalently $\approx 396 K^3$ Cirac--Zoller laser pulses). The paper spends
most of its pages building the arithmetic primitives (ADD, MULT, MODMULT,
EXP N) that reach these counts. Section~VII then closes the circle on the
smallest possible instance: factoring $N=15$.

For $N=15$ they exploit two special features. First, $x^4 \equiv 1 \pmod{15}$
for every $x$ coprime to 15, so the input register only needs $L=2$ qubits for
the modular-exponentiation step. Second, for $K=4$ and $L=2$ one can classically
precompute the entire 4-entry lookup table $a \mapsto x^a \bmod 15$ (only $2^L=4$
inputs) and then hand-craft a minimal quantum circuit that reproduces it.
For the worst-case bases $x=7$ or $x=13$ this gives their explicit Eq.~(7.5)
circuit on $L+K = 2+4 = 6$ qubits: 6 NOTs, 0 CNOTs, 4 Toffolis
$\Rightarrow$ 34 laser pulses. Adding the two Hadamards on the input
register (2 pulses) and the $L=2$ QFT (6 pulses) gives a full ``factor 15''
demonstration in \textbf{38 laser pulses} on the ion trap, via the further
optimized $\mathrm{EXP\,N}'$ variant (Eq.~7.9), which saves 4 pulses by
folding an input-superposition step into the exponentiation and using
Appendix A ``custom'' controlled-NOT gates.

\section{Claims table}

\begin{longtable}{@{}p{0.06\textwidth}p{0.52\textwidth}p{0.10\textwidth}p{0.10\textwidth}p{0.14\textwidth}@{}}
\toprule
\textbf{ID} & \textbf{Claim (from paper)} & \textbf{Testable classically?} & \textbf{Tested here?} & \textbf{Result} \\
\midrule
\endhead
C1 & $K$-bit number factorable in $\mathcal{O}(K^3)$ time on $5K+1$ qubits
     (Abstract, \S I). & Asymptotic; scaling only. & No (asymptotic). & N/A \\
C2 & Modular-exponential subcircuit costs $\approx 72 K^3$ elementary gates
     (Abstract, \S VI). & Analytic count. & No (asymptotic). & N/A \\
C3 & Ion-trap implementation of mod-exp costs $\approx 396 K^3$ laser pulses
     (Abstract, \S VI). & Analytic count. & No (asymptotic). & N/A \\
C4 & For $N=15$, general-purpose algorithm needs $L=2K=8$, 21 qubits,
     $\sim$15\,284 pulses (\S VII). & Concrete but big. & No (out of scope). & N/A \\
C5 & The specialized $N=15, L=2$ construction uses only $L+K=6$ qubits
     (\S VII). & Yes. & \textbf{Yes.} & \textbf{PASS} (6 qubits used) \\
C6 & Lookup table for $x=7$ is Eq.~(7.3): $\{0{\mapsto}1, 1{\mapsto}7,
     2{\mapsto}4, 3{\mapsto}13\}$. & Trivial classical. & \textbf{Yes.} &
     \textbf{PASS} (matches $7^a \bmod 15$). \\
C7 & The concrete circuit Eq.~(7.5) computes
     $|a\rangle|0\rangle \mapsto |a\rangle|7^a \bmod 15\rangle$. & Yes,
     small-instance sim. & \textbf{Yes.} & \textbf{PASS} for all 4 inputs
     with $p=1.000$. \\
C8 & On uniform superposition input, Eq.~(7.5) prepares Eq.~(7.2)
     $\frac{1}{2}\sum_a |a\rangle|7^a \bmod 15\rangle$. & Yes. & \textbf{Yes.} &
     \textbf{PASS} (fidelity $1.000000$). \\
C9 & Eq.~(7.5) has gate composition $[6, 0, 4]$ (Eq.~7.6). & Yes,
     literal count. & \textbf{Yes.} & \textbf{PASS} exactly. \\
C10 & Cirac--Zoller cost of Eq.~(7.5) is 34 laser pulses (\S VII text). &
     Yes, from pulse model. & \textbf{Yes.} & \textbf{PASS} (6$\times$1 +
     4$\times$7 = 34). \\
C11 & After a $L=2$ QFT on $\alpha$, $y$ is uniform on $\{0,1,2,3\}$
     (Eq.~7.10). & Yes, small-instance sim. & \textbf{Yes.} & \textbf{PASS}
     ($p(y{=}k) = 0.2500$ for each $k$). \\
C12 & Optimized $\mathrm{EXP\,N}'$ variant (Eq.~7.9) prepares Eq.~(7.2) in
     32 pulses; total 38 pulses to ``factor 15''. & Yes, but requires custom
     gate model. & Partial. & \textbf{SPOT-CHECK} (we identify the 4-pulse
     saving but do not re-simulate the custom gates; the arithmetic 32+6=38
     is verified). \\
\bottomrule
\end{longtable}

\section{Method}

All work was done inside a fresh Python 3.11 virtual environment on host
\texttt{CherryRd} (macOS 25.3.0).

\subsection{Tool versions}
\begin{itemize}
\item Python 3.11 (system interpreter)
\item \texttt{qiskit==2.5.0} (pypi)
\item \texttt{qiskit-aer==0.17.2} (pypi)
\item \texttt{numpy==2.4.6} (pypi)
\end{itemize}

\subsection{Exact commands}
\begin{lstlisting}[language=bash]
cd ~/Dropbox/REPLICATE-PROJECT/QC-200/QC-quant-ph-9602016-efficient-networks-quantum-factoring
curl -sL -o paper.pdf https://arxiv.org/pdf/quant-ph/9602016
pdftotext -layout paper.pdf work/paper.txt
python3.11 -m venv .venv
source .venv/bin/activate
pip install -q qiskit qiskit-aer numpy
python code/expn_7_15.py | tee logs/expn_7_15.log
\end{lstlisting}

\subsection{Circuit construction (verbatim from Eq.~7.5)}

We construct the operator as written in the paper, applying the rightmost
operator to the ket first (standard operator-product convention, consistent
with the ADD sequences of \S VI). The resulting Qiskit circuit is:

\begin{lstlisting}
                   +---+               +---+
alpha_0: -------o--| X |--o---------o--| X |--o-------
                |  +---+  |  +---+  |  +---+  |  +---+
alpha_1: -------o---------o--| X |--o---------o--| X |
         +---+  |       +-|-+ +---+  |         |  +---+
 beta_0: | X |--+-------| X |--------+---------+-------
         +---+  |       +---+        |       +-|-+
 beta_1: -------+---------------------+------| X |-----
         +---+  |                   +-|-+     +---+
 beta_2: | X |--+-------------------| X |-----------
         +---++-|-+                 +---+
 beta_3: -----| X |-----------------------
              +---+
count_ops = {'x': 6, 'ccx': 4}   depth = 8
\end{lstlisting}

\subsection{Verification protocol}

For each $a \in \{0,1,2,3\}$ we (i) prepare $|a\rangle_\alpha |0\rangle_\beta$
by X-ing the appropriate $\alpha$ qubits, (ii) apply the Eq.~(7.5) circuit,
(iii) compute the full 6-qubit statevector, (iv) read the single non-zero
computational basis state (probability $= 1.000$ in each case, confirming
the operator is a permutation on the basis), and (v) decode $(a_\mathrm{out},
b_\mathrm{out})$ from the bitstring. We then check
$b_\mathrm{out} = 7^{a_\mathrm{in}} \bmod 15$.

For the entangled-state verification we prepare
$(1/2)\sum_a |a\rangle |0\rangle$ via two Hadamards on $\alpha$, apply
Eq.~(7.5), and take the fidelity against the analytical target Eq.~(7.2)
constructed by direct amplitude assignment.

For the period-finding demo we apply Qiskit's \texttt{QFT(2, do\_swaps=True)}
to the $\alpha$ register and marginalize the full statevector over $\beta$
to obtain the $y$ distribution.

\section{Results (vs.\ paper)}

\begin{center}
\begin{tabular}{@{}lllr@{}}
\toprule
\textbf{Quantity} & \textbf{Paper value} & \textbf{Our value} & \textbf{Match} \\
\midrule
Lookup table $x=7$ (Eq.~7.3)      & \{1, 7, 4, 13\} & \{1, 7, 4, 13\} & PASS \\
Per-input basis-state probability & 1               & 1.000           & PASS \\
Entangled-state fidelity          & 1               & 1.000000        & PASS \\
NOT count in Eq.~(7.5)            & 6               & 6               & PASS \\
CNOT count in Eq.~(7.5)           & 0               & 0               & PASS \\
Toffoli count in Eq.~(7.5)        & 4               & 4               & PASS \\
Cirac--Zoller pulses for EXP N    & 34              & 34              & PASS \\
$p(y)$ after $L=2$ QFT (each $y$) & $1/4 = 0.2500$  & 0.2500          & PASS \\
Total pulses (first construction) & 42              & 42              & PASS \\
Total pulses (optimized EXP N$'$) & 38              & --- (not re-sim.) & N/A \\
\bottomrule
\end{tabular}
\end{center}

Full machine-readable results are in
\path{report/evidence/expn_7_15_result.json}.

\section{Verdict}

\textbf{REPLICATED} for the Sec.~VII $N=15$ headline construction.

\emph{Justification.} We reproduced the paper's explicit
$\mathrm{EXP\,N}(x{=}7, N{=}15)$ operator (Eq.~7.5) exactly as written and
verified by direct Qiskit statevector simulation that (a) it computes the
Eq.~(7.3) lookup table on all 4 inputs with unit probability, (b) it prepares
the entangled state of Eq.~(7.2) on 6 qubits with fidelity 1, (c) its gate
composition is exactly $[6, 0, 4]$ (Eq.~7.6), (d) its Cirac--Zoller pulse cost
is exactly 34 (per the pulse model of \S III), and (e) the QFT-based
period-finding step gives the predicted uniform distribution of $y$
(Eq.~7.10). The abstract's ``38 pulses'' headline requires the further
optimized $\mathrm{EXP\,N}'$ variant (Eq.~7.9) using Appendix A custom gates;
we identify it and verify the arithmetic $32 + 6 = 38$, but do not
independently pulse-count the custom controlled-NOTs (which requires
implementing the Appendix A gate decompositions, out of scope for this
short replication).

\section*{Open Questions}

\begin{enumerate}
  \item[Q1.] \textbf{Custom-gate implementation cost.} The paper claims each
    ``custom'' $C_{[[i,j]],k}$ gate (Appendix A) costs 7 laser pulses, the
    same as a standard Toffoli, yet enables the 4-pulse saving from
    $\mathrm{EXP\,N}$ (Eq.~7.5) to $\mathrm{EXP\,N}'$ (Eq.~7.9). Is that pulse
    count really identical, or does it hide a phase-tracking overhead that a
    fair modern noise model would surface? A direct re-implementation on a
    Cirac--Zoller pulse-level simulator (e.g.\ \texttt{QuTiP}) would settle
    this.
  \item[Q2.] \textbf{Noise-inclusive break-even.} The paper is purely
    coherent: it counts pulses but never asks at what per-pulse depolarizing
    error rate $p$ the full 38-pulse ``factor 15'' demo actually succeeds
    with probability $\ge 1/2$. Given a plausible ion-trap $p \sim
    10^{-3}$--$10^{-2}$ of the era, what is the crossover, and does the
    4-pulse $\mathrm{EXP\,N}'$ optimization actually matter under noise?
  \item[Q3.] \textbf{Ordering ambiguity of Eq.~(7.5).} Reading Eq.~(7.5)
    left-to-right vs.\ right-to-left changes which basis state each Toffoli
    corrects; we found that right-to-left (rightmost applied first) is
    consistent with the paper's ADD conventions in \S VI (Eq.~6.36) and
    reproduces the lookup table. But the paper never states this
    convention explicitly. Is there a sub-corpus of Preskill-era papers where
    left-to-right is used, and how many replications of pre-2000
    ``operator-product'' circuits silently assume the wrong direction?
  \item[Q4.] \textbf{Beyond $N=15$: how quickly does the special-case trick
    stop working?} The Eq.~(7.5) construction exploits $r(x) \mid 4$ for all
    $x \in (\mathbb{Z}/15\mathbb{Z})^*$. For $N = 21, 33, 35$, the orders
    quickly exceed $2^L$ for any small $L$, forcing the general-purpose
    algorithm. What is the smallest $N$ for which no analogous hand-crafted
    single-Toffoli-layer $\mathrm{EXP\,N}$ exists, and does its custom
    minimal-pulse circuit still beat the general $72 K^3$ construction?
  \item[Q5.] \textbf{Reproducibility of the $\sim$15\,000-pulse general-purpose
    estimate for $N=15$.} Section~VII quotes 15\,284 pulses on 21 qubits and
    14\,878 on 22, but the paper does not release code. A modern
    re-derivation running the full ADD/MULT/MODMULT/EXPN cascade at
    $K=4, L=8$ would either confirm or refute those numbers, and would be a
    natural next step for this replication: it is finite, small enough to
    simulate exhaustively, and would validate the general-purpose complexity
    bounds that underpin the abstract's $72 K^3$ headline number.
\end{enumerate}

\end{document}
